\documentclass[11pt,a4paper]{article}

% Packages
\usepackage[utf8]{inputenc}
\usepackage[T1]{fontenc}
\usepackage{geometry}
\usepackage{graphicx}
\usepackage{booktabs}
\usepackage{hyperref}
\usepackage{xcolor}
\usepackage{amsmath}
\usepackage{amssymb}
\usepackage{enumitem}
\usepackage{fancyhdr}
\usepackage{titlesec}
\usepackage{tcolorbox}
\usepackage{array}
\usepackage{multirow}

% Page geometry
\geometry{margin=1in}

% Colors
\definecolor{sectionblue}{RGB}{0, 83, 156}
\definecolor{warningorange}{RGB}{255, 152, 0}
\definecolor{successgreen}{RGB}{76, 175, 80}
\definecolor{lightgray}{RGB}{245, 245, 245}

% Section formatting
\titleformat{\section}{\Large\bfseries\color{sectionblue}}{\thesection}{1em}{}
\titleformat{\subsection}{\large\bfseries}{\thesubsection}{1em}{}

% Header/Footer
\pagestyle{fancy}
\fancyhf{}
\rhead{Model Card: ResNet-GPT2 Maze Solver}
\lhead{Vision-to-Language Navigation}
\rfoot{Page \thepage}

% Custom boxes
\newtcolorbox{infobox}[1][]{colback=lightgray,colframe=sectionblue,title=#1,fonttitle=\bfseries}
\newtcolorbox{warningbox}{colback=warningorange!10,colframe=warningorange,title=\textbf{Limitations \& Risks}}
\newtcolorbox{metricbox}{colback=successgreen!10,colframe=successgreen}

\begin{document}

% Title
\begin{center}
    {\Huge\bfseries\color{sectionblue} Model Card}\\[0.5cm]
    {\LARGE ResNet18-GPT2 Vision-to-Language Maze Solver}\\[0.3cm]
    {\large A Prefix-Tuning Approach for Visual Navigation Sequence Generation}\\[0.5cm]
    \rule{\textwidth}{1pt}\\[0.3cm]
    {\large Seena Mohajeran and Arshdeep Dhillon}
\end{center}

\vspace{0.5cm}

%==============================================================================
\section{Model Overview}
%==============================================================================

\begin{infobox}[Quick Reference]
\begin{tabular}{@{}ll}
\textbf{Model Type:} & Vision-to-Language (Image Captioning variant) \\
\textbf{Architecture:} & ResNet18 encoder + GPT-2 decoder with prefix-tuning \\
\textbf{Task:} & Generate navigation sequences from maze images \\
\textbf{Input:} & RGB maze image (224$\times$224 pixels) \\
\textbf{Output:} & Sequence of directional tokens (R, U) \\
\textbf{Parameters:} & 12,215,488 total \\
\end{tabular}
\end{infobox}

\subsection{Model Description}

This model investigates whether visual spatial information can be effectively encoded into language model representations through prefix-tuning. The architecture combines a pretrained ResNet18 vision encoder with a small GPT-2 language decoder, connected via learned prefix embeddings that project visual features into the language model's embedding space.

The model takes a maze image as input and autoregressively generates a sequence of navigation commands (Right ``R'' and Up ``U'') that traverse from the start position (bottom-left) to the goal position (top-right). This task requires the model to:
\begin{enumerate}[noitemsep]
    \item Extract spatial structure from the visual input
    \item Encode wall positions and valid paths
    \item Generate a coherent sequence respecting maze constraints
\end{enumerate}

\subsection{Intended Use}

\textbf{Primary Use Cases:}
\begin{itemize}[noitemsep]
    \item Research into vision-to-language information transfer
    \item Investigating prefix-tuning as a cross-modal bridging mechanism
    \item Educational demonstration of multimodal architectures
    \item Baseline for visual navigation tasks
\end{itemize}

\textbf{Out-of-Scope Uses:}
\begin{itemize}[noitemsep]
    \item Real-world robot navigation
    \item Safety-critical pathfinding applications
    \item General-purpose image captioning
    \item Mazes with different start/goal configurations
\end{itemize}

%==============================================================================
\section{Architecture Details}
%==============================================================================

\subsection{Model Architecture}

\begin{center}
\begin{tabular}{lll}
\toprule
\textbf{Component} & \textbf{Specification} & \textbf{Details} \\
\midrule
\multirow{3}{*}{Vision Encoder} & Base Model & ResNet18 (ImageNet pretrained) \\
& Output Dimension & 512 \\
& Training Status & Frozen \\
\midrule
\multirow{2}{*}{Feature Projection} & Architecture & Linear(512 $\rightarrow$ 128) \\
& Activation & None \\
\midrule
\multirow{3}{*}{Prefix Generator} & Layer 1 & Linear(128 $\rightarrow$ 256) + Tanh \\
& Layer 2 & Linear(256 $\rightarrow$ 16 $\times$ 128) \\
& Output Shape & [batch, 16, 128] \\
\midrule
\multirow{5}{*}{GPT-2 Decoder} & Hidden Size & 128 \\
& Layers & 2 \\
& Attention Heads & 2 \\
& Max Positions & 128 \\
& Dropout & 0.4 \\
\bottomrule
\end{tabular}
\end{center}

\subsection{Tokenizer}

\begin{center}
\begin{tabular}{cl}
\toprule
\textbf{Token ID} & \textbf{Token} \\
\midrule
0 & \texttt{<pad>} (Padding) \\
1 & \texttt{<s>} (Beginning of Sequence) \\
2 & \texttt{</s>} (End of Sequence) \\
3 & \texttt{<unk>} (Unknown) \\
4 & \texttt{R} (Move Right) \\
5 & \texttt{U} (Move Up) \\
\bottomrule
\end{tabular}
\end{center}

\textbf{Vocabulary Size:} 6 tokens

%==============================================================================
\section{Training Details}
%==============================================================================

\subsection{Training Data}

\begin{center}
\begin{tabular}{ll}
\toprule
\textbf{Parameter} & \textbf{Value} \\
\midrule
Grid Size & 7$\times$7 \\
Training Samples & 36,950 \\
Test Samples & 9,250 \\
Variations per Solution & 50 \\
Start Position & (0, 6) --- Bottom-left \\
Goal Position & (6, 0) --- Top-right \\
Random Seed & 69 \\
Image Resolution & 224$\times$224 RGB \\
\bottomrule
\end{tabular}
\end{center}

The dataset consists of procedurally generated maze images with corresponding solution paths. Each unique maze solution is augmented with 50 visual variations to improve generalization. Mazes use a binary representation where white cells indicate valid paths and black cells indicate walls.

\subsection{Training Configuration}

\begin{center}
\begin{tabular}{ll}
\toprule
\textbf{Hyperparameter} & \textbf{Value} \\
\midrule
Epochs & 100 \\
Batch Size & 32 \\
Base Learning Rate & $5 \times 10^{-4}$ \\
Optimizer & AdamW \\
Weight Decay & 0.01 \\
LR Schedule & Cosine Annealing \\
Minimum LR & $1 \times 10^{-6}$ \\
Gradient Clipping & Max norm 1.0 \\
\bottomrule
\end{tabular}
\end{center}

\textbf{Differentiated Learning Rates:}
\begin{itemize}[noitemsep]
    \item ResNet features: $\text{lr} / 10$ (frozen, effectively 0)
    \item Feature projection: $\text{lr}$
    \item Prefix generator: $\text{lr}$
    \item GPT-2 decoder: $\text{lr} / 5$
\end{itemize}

\subsection{Image Preprocessing}

\begin{enumerate}[noitemsep]
    \item Resize to 224$\times$224 pixels
    \item Convert to tensor (scale to [0, 1])
    \item Normalize with ImageNet statistics:
    \begin{itemize}[noitemsep]
        \item Mean: [0.485, 0.456, 0.406]
        \item Std: [0.229, 0.224, 0.225]
    \end{itemize}
\end{enumerate}

%==============================================================================
\section{Limitations}
%==============================================================================

\begin{warningbox}
\begin{enumerate}[noitemsep]
    \item \textbf{Grid Size Constraint:} Trained exclusively on 7$\times$7 grids; performance on other sizes is untested and likely poor.
    
    \item \textbf{Fixed Start/Goal:} Only supports bottom-left to top-right navigation; cannot handle arbitrary start/goal positions.
    
    \item \textbf{Limited Action Space:} Only Right (R) and Up (U) moves; cannot handle mazes requiring Left or Down movements.
    
    \item \textbf{Capacity Bottleneck:} The 128-dimensional hidden size and 16 prefix tokens may be insufficient for encoding complex maze structures. Prior experiments on smaller grids suggest this is a key limiting factor.
    
    \item \textbf{Wall Collision Errors:} The majority of failures (99.3\%) involve hitting walls, indicating imprecise spatial encoding.
    
    \item \textbf{No Uncertainty Quantification:} Model provides no confidence scores or indication of prediction reliability.
    
    \item \textbf{Single Solution Generation:} Cannot generate multiple alternative paths or rank solutions.
\end{enumerate}
\end{warningbox}

%==============================================================================
\section{Technical Specifications}
%==============================================================================

\subsection{Compute Infrastructure}

\begin{center}
\begin{tabular}{ll}
\toprule
\textbf{Component} & \textbf{Specification} \\
\midrule
Hardware & CUDA-compatible GPU \\
Framework & PyTorch \\
Vision Backbone & torchvision (ResNet18) \\
Language Model & Hugging Face Transformers (GPT2LMHeadModel) \\
Python Version & 3.12 \\
\bottomrule
\end{tabular}
\end{center}

\subsection{Model Files}

\begin{center}
\begin{tabular}{ll}
\toprule
\textbf{File} & \textbf{Description} \\
\midrule
\texttt{model.py} & Model architecture definition \\
\texttt{tokenizer.py} & Sequence tokenizer \\
\texttt{dataset.py} & Dataset and DataLoader utilities \\
\texttt{train\_utils.py} & Training loop and evaluation \\
\texttt{solution\_validator.py} & Path validation logic \\
\texttt{data\_utils.py} & Data loading utilities \\
\texttt{resnet\_gpt2\_prefix.pth} & Trained model checkpoint \\
\bottomrule
\end{tabular}
\end{center}

\subsection{Inference Example}

\begin{verbatim}
from model import ResNetGPT2PrefixModel
from tokenizer import SimpleTokenizer

# Load model
model = ResNetGPT2PrefixModel(
    vocab_size=6,
    hidden_size=128,
    num_layers=2,
    num_attention_heads=2,
    num_prefix_tokens=16
)
model.load_state_dict(checkpoint['model_state_dict'])

# Generate solution
output = model.generate(
    images,           # [batch, 3, 224, 224]
    max_length=20,
    bos_token_id=1,
    eos_token_id=2,
    pad_token_id=0
)
\end{verbatim}

%==============================================================================
\section{Citation}
%==============================================================================

If you use this model or methodology in your research, please cite:

\begin{verbatim}
@misc{resnet_gpt2_maze_solver,
  title={ResNet-GPT2 Vision-to-Language Maze Solver: 
         A Prefix-Tuning Approach},
  author={[Author Name]},
  year={2025},
  note={Model Card v1.0}
}
\end{verbatim}

\vspace{1cm}
\begin{center}
\rule{0.5\textwidth}{0.5pt}\\[0.3cm]
{\small\textit{This model card follows the format proposed by Mitchell et al. (2019) ``Model Cards for Model Reporting''}}
\end{center}

\end{document}
